\chapter{Types d'ordres et leurs usages}\label{ch:mx:order-types-and-their-uses}

En janvier 2015, deux bourses américaines ont accepté de verser une pénalité de 14~millions de dollars, sans admettre ni nier les constats, pour régler des poursuites de la SEC
portant sur leurs types d'ordres. Leurs règlements décrivaient une seule façon de réajuster le prix d'un ordre qui aurait bloqué ou croisé la cotation d'une autre plateforme~;
les bourses en acceptaient trois, avec des priorités différentes entre elles, et n'en avaient décrit complètement une qu'à certains membres, mais pas à tous. Rien n'était
caché dans le code du \omterm{def:mx:the-limit-order-book:engine}{moteur d'appariement}. Il était caché dans l'écart entre ce que le moteur faisait d'un ordre et ce que le règlement de marché annonçait
qu'il ferait. Un type d'ordre est exactement cela~: un contrat sur ce que le moteur fera d'un ordre, et ce sont les petits caractères qui décident de qui exécute en premier.

\section{Conditions de prix~: cours limité, au marché, cours limité exécutable, stop}

Un \omterm{def:mx:the-limit-order-book:orders}{ordre à cours limité} et un \omterm{def:mx:the-limit-order-book:orders}{ordre au marché} (\cref{def:mx:the-limit-order-book:orders}) diffèrent par ce qu'ils garantissent~: l'\omterm{def:mx:the-limit-order-book:orders}{ordre à cours limité} garantit son prix, l'\omterm{def:mx:the-limit-order-book:orders}{ordre au marché}
son exécution, tant qu'il y a quelque chose à exécuter contre lui. Entre les deux se situe l'ordre que les intervenants envoient le plus souvent.

\begin{definition}[Ordre à cours limité exécutable]\label{def:mx:order-types-and-their-uses:marketable}
Un \emph{ordre à cours limité exécutable}\index{ordre à cours limité exécutable} est un \omterm{def:mx:the-limit-order-book:orders}{ordre à cours limité} dont le prix atteint, à son arrivée, le meilleur cours affiché en face~: un achat à un cours égal ou
supérieur au meilleur cours vendeur, une vente à un cours égal ou inférieur au meilleur cours acheteur. Il s'exécute immédiatement contre le carnet, jusqu'à son cours limite, et le reliquat éventuel reste affiché à ce prix.
\end{definition}

Un \omterm{def:mx:order-types-and-their-uses:marketable}{ordre à cours limité exécutable} est un \omterm{def:mx:the-limit-order-book:orders}{ordre au marché} muni d'un garde-fou~: un acheteur qui envoie 200 actions au cours vendeur, majoré de cinq centimes, exécute tout ce qui est disponible
jusqu'à ce prix et rien au-delà. Ce garde-fou compte lorsque le carnet est peu profond, et il compte surtout lorsque les ordres ne sont pas envoyés par une personne qui réagit au marché
mais par un prix que le marché a atteint.

\begin{definition}[Ordre stop, ordre stop à cours limité]\label{def:mx:order-types-and-their-uses:stop}
Un \emph{ordre stop}\index{ordre stop} reste invisible jusqu'à ce qu'une transaction ait lieu à son cours de déclenchement ou au-delà (à ce cours ou au-dessus pour un stop d'achat, à ce cours ou en dessous pour
un stop de vente)~; il devient alors un \omterm{def:mx:the-limit-order-book:orders}{ordre au marché}. Un \emph{ordre stop à cours limité}\index{ordre stop à cours limité} devient à la place un \omterm{def:mx:the-limit-order-book:orders}{ordre à cours limité}, à son cours limite déclaré.
\end{definition}

Les stops sont des ordres de ceux qui veulent sortir lorsqu'un cours est atteint~: c'est le stop-loss. Leurs déclencheurs sont des prix, leur exécution est immédiate, et
un carnet rempli de stops est un carnet qui vend d'autant plus fort qu'il baisse. La \Cref{fig:mx:order-types-and-their-uses:cascade} mesure cet effet sur
\texttt{firm.exchsim}, le simulateur de bourse du livre (construit au chapitre~26)~: une échelle d'achats de 500 actions à chaque centime en dessous de 100.00, des stops de vente de 800 actions
à chaque centime à partir de 99.95 vers le bas, et une vente au marché de 3\,000 actions. Sans stops, la vente descend l'échelle de cinq centimes. Avec vingt stops en dessous,
chaque vente déclenchée absorbe davantage d'achats, atteint le déclencheur suivant, et le cours chute de 37 centimes~; 19\,000 actions s'échangent pour une vente initiale de 3\,000.
Les \omterm{def:mx:order-types-and-their-uses:stop}{ordres stop à cours limité}, fixés à un centime sous leur déclencheur, arrêtent la cascade plus tôt (25 centimes), car un ordre qui ne peut pas vendre à son cours limite reste affiché
au lieu de descendre l'échelle~: le garde-fou, encore, payé au prix de ne pas sortir.

\begin{omfigure}
\begin{tikzpicture}
\begin{axis}[width=8.6cm,height=5.2cm,xlabel={\small stops de vente sous le marché},ylabel={\small baisse du cours (centimes)},
  tick label style={font=\footnotesize},xtick={0,2,4,6,8,10,12,14,16,18,20},ymin=0,ymax=40,grid=major,grid style={gray!25},
  legend columns=2,legend style={font=\footnotesize,at={(0.5,-0.3)},anchor=north,draw=none,fill=none,
  /tikz/every even column/.append style={column sep=8pt}},table/col sep=comma]
\addplot[omThm,very thick,mark=*] table[x=n,y=stop_market]{figdata/microstructure/02-order-types-and-their-uses/cascade.csv};
\addplot[omDef,very thick,mark=square*] table[x=n,y=stop_limit]{figdata/microstructure/02-order-types-and-their-uses/cascade.csv};
\legend{ordres stop,ordres stop à cours limité (un centime sous le déclencheur)}
\end{axis}
\end{tikzpicture}

\omcaption{Une cascade de stops sur \texttt{firm.exchsim}~: la baisse provoquée par une vente au marché de 3\,000 actions dans une échelle de 500 actions par centime, en regard du
nombre de stops de vente de 800 actions placés à un centime d'écart à partir de 99.95 vers le bas. Données~: \texttt{mx\_ordertypes.stop\_cascade}.}\label{fig:mx:order-types-and-their-uses:cascade}
\end{omfigure}

\section{Conditions de durée}

\begin{definition}[Durée de validité~; ordres exécuter-ou-annuler, tout-ou-rien et à révocation]\label{def:mx:order-types-and-their-uses:tif}
La \emph{durée de validité}\index{durée de validité} (time in force) d'un ordre indique combien de temps sa partie non exécutée peut rester affichée. Un ordre du jour reste affiché jusqu'à la fin de la séance. Un
\emph{ordre exécuter-ou-annuler}\index{ordre exécuter-ou-annuler} (immediate-or-cancel, IOC) s'exécute dans la mesure du possible à son arrivée et annule le reste. Un \emph{ordre tout-ou-rien
}\index{ordre tout-ou-rien} (fill-or-kill, FOK) s'exécute en totalité à son arrivée, ou pas du tout. Un \emph{ordre à révocation}\index{ordre à révocation} (good-till-cancelled, GTC)
survit à la fin de la séance et reste affiché jusqu'à ce qu'il soit exécuté ou annulé.
\end{definition}

Les ordres IOC et FOK ne restent jamais affichés, si bien qu'ils ne montrent rien au marché avant d'être négociés~: ce sont les moyens qu'utilisent les algorithmes pour prélever de la liquidité, balayer plusieurs
plateformes à la fois (l'ordre de balayage intermarchés du livre~1, chapitre~9, est un IOC muni d'un indicateur juridique) et sonder les dark pools (chapitre~9). Les \omterm{def:mx:order-types-and-their-uses:tif}{durées de validité} à l'ouverture
et à la clôture, qui ne restent affichées que pour une enchère, relèvent du chapitre~10. Le protocole de saisie des ordres du simulateur porte les six
(\texttt{D}, \texttt{G}, \texttt{I}, \texttt{F}, \texttt{O}, \texttt{C}).

Le tableau montre ce qu'un même carnet fait de huit ordres de natures différentes. La firme~1 affiche 300 actions à l'achat à 99.99 et 500 à la vente à 100.01~; la firme~2 envoie ensuite,
dans l'ordre~:

\begin{center}
\small
\begin{tabular}{@{}p{4.4cm}p{3.2cm}rp{4.6cm}@{}}
\toprule
ordre & réponse & exécuté & ce qui s'est passé \\
\midrule
achat à cours limité de 100 à 100.00 & accepté & 0 & affiché~; devient le meilleur cours acheteur, 100 affichées \\
achat à cours limité exécutable de 200 à 100.01 & accepté, exécuté & 200 & prend 200 des 500 à la vente \\
achat post-only de 100 à 100.01 & rejeté (croiserait) & 0 & n'entre jamais dans le carnet \\
achat tout-ou-rien de 400 à 100.01 & accepté, annulé & 0 & seules 300 restent à 100.01~: rien ne s'échange \\
achat exécuter-ou-annuler de 400 à 100.01 & accepté, exécuté, annulé & 300 & prend les 300, annule 100 \\
vente iceberg de 1\,000 à 100.02, 100 affichées & accepté & 0 & affiché~; le flux n'en montre que 100 \\
vente indexée au mid de 200 & accepté & 0 & reste caché au mid, à 100.01 \\
vente stop de 300 à 99.99 & accepté (en attente) & 0 & invisible jusqu'à une transaction à 99.99 ou en dessous \\
\bottomrule
\end{tabular}
\end{center}

\section{Conditions d'affichage~: caché, iceberg, indexé}

\begin{definition}[Ordre caché, ordre iceberg]\label{def:mx:order-types-and-their-uses:hidden}
Un \emph{ordre caché}\index{ordre caché} reste dans le carnet et peut s'exécuter contre les ordres entrants, mais il n'apparaît dans aucun flux de données de marché. Un \emph{ordre iceberg
}\index{ordre iceberg} (ordre de réserve) n'affiche qu'une tranche de sa taille~; lorsque cette tranche est exécutée, la plateforme affiche la tranche suivante prélevée sur la
réserve, à la fin de la file.
\end{definition}

Un \omterm{def:mx:order-types-and-their-uses:hidden}{ordre caché} renonce à deux choses pour rester secret. À son cours, il passe après tous les ordres affichés, quelles que soient leurs heures d'arrivée (la priorité
qu'implémente \texttt{firm.lob}), et, étant invisible, il n'attire personne~: un cours acheteur affiché indique aux vendeurs où vendre, tandis qu'un \omterm{def:mx:order-types-and-their-uses:hidden}{ordre caché} attend d'être trouvé.
La tranche réapprovisionnée d'un iceberg est une nouvelle arrivée et rejoint la fin de la file avec une nouvelle référence dans le flux (chapitre~1), si bien qu'un grand iceberg
n'est négocié qu'une tranche à chaque passage dans la file.

\begin{definition}[Ordre indexé, ordre indexé au mid]\label{def:mx:order-types-and-their-uses:peg}
Un \emph{ordre indexé}\index{ordre indexé} (pegged order) n'a pas de prix fixe~: la plateforme le fixe à partir d'une cotation de référence et le réajuste chaque fois que la référence bouge,
dans une limite facultative. Un ordre indexé primaire suit le meilleur cours de son propre côté~; un \emph{ordre indexé au mid}\index{ordre indexé au mid} (midpoint peg) reste caché au milieu
du meilleur cours acheteur et du meilleur cours vendeur.
\end{definition}

Un \omterm{def:mx:order-types-and-their-uses:peg}{ordre indexé au mid} s'exécute à mi-fourchette~: un acheteur et un vendeur qui s'indexent tous deux au milieu s'y retrouvent, chacun économisant une demi-fourchette plutôt que de la traverser.
C'est le métier des dark pools (chapitre~9), et la raison pour laquelle les bourses proposent aussi des \omterm{def:mx:order-types-and-their-uses:peg}{ordres indexés au mid}. Le texte de règlement d'une bourse montre à quel point
ce type d'ordre porte de petits caractères.

\begin{dated}{2026-06}{Indexation et quantité minimale sur Nasdaq}\label{dat:mx:order-types-and-their-uses:nasdaq}
La Rule 4703 de Nasdaq, telle que reproduite dans une demande de modification de règlement de juin 2026~: un ordre avec Pegging reçoit un nouvel horodatage chaque fois que son cours est mis à jour (il est
évalué comme un ordre nouvellement saisi)~; un ordre avec Midpoint Pegging est retiré lorsque le meilleur cours acheteur et le meilleur cours vendeur internes deviennent croisés ou disparaissent, est ressaisi au
nouveau milieu, et annulé si aucun cours valide ne revient dans un délai d'une seconde~; les \omterm{def:mx:order-types-and-their-uses:peg}{ordres indexés au mid} sont annulés lorsqu'une suspension de cotation est déclarée~;
les \omterm{def:mx:order-types-and-their-uses:peg}{ordres indexés} sont soumis à un collar et sont annulés lorsqu'ils s'exécuteraient à plus de 0.25~\$ ou 5~\% (selon le plus élevé des deux) de moins bien que les
meilleurs cours acheteur et vendeur nationaux~; un ordre muni de l'attribut Minimum Quantity doit être au moins égal à un lot standard et ne peut pas être affiché.
\end{dated}

\begin{omfigure}
\begin{tikzpicture}[font=\footnotesize,x=1cm,y=1cm]
\node[anchor=west,font=\small\bfseries] at (-0.2,3.3) {iceberg};
\foreach \x/\w/\l in {0/1.3/A,1.35/0.9/B} {\draw[fill=omDef!18,draw=omDef] (\x,2.4) rectangle ++(\w,0.55); \node at (\x+\w/2,2.675) {\l};}
\draw[fill=omExo!40,draw=omExo,thick] (2.3,2.4) rectangle ++(0.95,0.55); \node[font=\scriptsize] at (2.775,2.675) {tranche};
\draw[fill=gray!12,draw=gray,dashed] (2.3,1.6) rectangle ++(2.5,0.55); \node at (3.55,1.875) {réserve cachée};
\draw[->,thick,omExo] (3.3,2.675) to[out=0,in=180] (5.0,2.675);
\draw[fill=omExo!40,draw=omExo,thick] (5.05,2.4) rectangle ++(0.95,0.55); \node[font=\scriptsize] at (5.525,2.675) {suivante};
\node[anchor=west,align=left] at (6.1,2.675) {lorsque la tranche est exécutée,\\ la suivante rejoint la fin de la file,\\ sous une nouvelle référence};
\node[anchor=west,font=\small\bfseries] at (-0.2,1.0) {ordre indexé au mid};
\draw[thick,omDef] (0,0) -- (3,0); \node[anchor=east] at (-0.05,0) {cours acheteur 100.00};
\draw[thick,omThm] (0,0.5) -- (3,0.5); \node[anchor=east] at (-0.05,0.5) {cours vendeur 100.02};
\draw[dashed,gray] (0,0.25) -- (3,0.25);
\draw[fill=gray!20,draw=gray] (1.1,0.13) rectangle ++(0.8,0.24);
\node[anchor=west,align=left] at (3.3,0.25) {reste caché au mid, 100.01~; se déplace dès que\\ l'un des deux cours affichés bouge~; s'échange\\ à mi-fourchette contre l'un ou l'autre côté};
\end{tikzpicture}

\omcaption{Deux conditions d'affichage. En haut~: un iceberg n'affiche qu'une tranche~; lorsqu'elle est consommée, la plateforme affiche la tranche suivante en queue de file.
En bas~: un \omterm{def:mx:order-types-and-their-uses:peg}{ordre indexé au mid} reste, invisible, à mi-chemin entre le meilleur cours acheteur et le meilleur cours vendeur.}\label{fig:mx:order-types-and-their-uses:display}
\end{omfigure}

Un iceberg cache sa taille mais non son comportement. Sur \texttt{firm.exchsim}, avec le flux simulé du livre~7 en toile de fond, un agent vend 30\,000 actions
au meilleur cours vendeur, en n'en affichant que 300 à la fois, d'abord sous forme d'iceberg de plateforme, puis sous forme d'algorithme qui envoie lui-même un nouvel ordre de 300 actions après chaque exécution de tranche. Un
détecteur ne lit que le flux public~: après qu'une exécution a retiré entièrement un ordre affiché, il cherche un ajout du même côté et au même prix dans les 50
millisecondes et de taille exactement égale à celle de l'ordre retiré. La \Cref{fig:mx:order-types-and-their-uses:iceberg} montre la course aux armements. Le réapprovisionnement de la plateforme arrive
dans le même événement que l'exécution et il est repéré à chaque fois (99 sur 99), avec 20 fausses alertes sur la demi-heure, dues à d'autres ordres qui par hasard
correspondaient. Un algorithme qui se réapprovisionne aussitôt est repéré dans 62~\% des cas (ses réapprovisionnements qui ont rencontré un marché ayant bougé ont été exécutés dès leur arrivée et sont restés
au carnet, de taille réduite)~; un réapprovisionnement après un délai aléatoire de 200 millisecondes en moyenne ramène ce taux à 18~\%, et une taille tirée au hasard, à 30~\% près, le ramène à 6~\%. Assouplir le
détecteur (une seconde, tailles à 35~\% près) ramène la détection à 47~\%, mais avec 378 fausses alertes pour 34 détections. Ce sont les délais et les tailles aléatoires qui
expliquent pourquoi les algorithmes d'exécution qui réapprovisionnent leurs propres ordres sont plus difficiles à repérer que les icebergs de plateforme.

\begin{omfigure}
\begin{tikzpicture}
\begin{axis}[name=rate,width=7.2cm,height=5.2cm,title={\small réapprovisionnements détectés},ybar,bar width=7pt,ymin=0,ymax=1.1,
  xtick={0,1,2,3},xticklabels={plateforme,algo,algo 200 ms,algo 200 ms $\pm$30~\%},x tick label style={font=\footnotesize,rotate=20,anchor=east},
  tick label style={font=\footnotesize},grid=major,grid style={gray!25},enlarge x limits=0.18,
  legend columns=2,legend style={font=\footnotesize,at={(1.12,-0.42)},anchor=north,draw=none,fill=none,
  /tikz/every even column/.append style={column sep=8pt}},table/col sep=comma]
\addplot[fill=omDef!60,draw=omDef] table[x=k,y=strict]{figdata/microstructure/02-order-types-and-their-uses/iceberg.csv};
\addplot[fill=omThm!60,draw=omThm] table[x=k,y=loose]{figdata/microstructure/02-order-types-and-their-uses/iceberg.csv};
\legend{détecteur strict (50 ms~; taille exacte),détecteur souple (1 s~; taille à 35~\% près)}
\end{axis}
\begin{axis}[at={(rate.east)},anchor=west,xshift=1.2cm,width=7.2cm,height=5.2cm,title={\small fausses alertes en une demi-heure},ybar,bar width=7pt,
  ymin=0,ymax=420,xtick={0,1,2,3},xticklabels={plateforme,algo,algo 200 ms,algo 200 ms $\pm$30~\%},
  x tick label style={font=\footnotesize,rotate=20,anchor=east},tick label style={font=\footnotesize},grid=major,grid style={gray!25},
  enlarge x limits=0.18,table/col sep=comma]
\addplot[fill=omDef!60,draw=omDef] table[x=k,y=strict_fa]{figdata/microstructure/02-order-types-and-their-uses/iceberg.csv};
\addplot[fill=omThm!60,draw=omThm] table[x=k,y=loose_fa]{figdata/microstructure/02-order-types-and-their-uses/iceberg.csv};
\end{axis}
\end{tikzpicture}

\omcaption{Détection des icebergs à partir du flux public~: la part des vrais réapprovisionnements que signalent un détecteur strict et un détecteur souple (à gauche), ainsi que leurs fausses alertes (à droite),
pour un iceberg de plateforme et pour un algorithme qui se réapprovisionne aussitôt, après un délai aléatoire (moyenne de 200~ms), et avec une taille aléatoire (jusqu'à 30~\% dans un sens ou l'autre)~;
30\,000 actions affichées 300 à la fois sur une demi-heure simulée. Données~: \texttt{mx\_ordertypes.iceberg\_study}.}\label{fig:mx:order-types-and-their-uses:iceberg}
\end{omfigure}

\section{Qui utilise quoi, et pourquoi}

\begin{center}
\small
\begin{tabular}{@{}p{3.1cm}p{4.9cm}p{5.6cm}@{}}
\toprule
participant & ordres typiques & pourquoi \\
\midrule
investisseur particulier & au marché, à cours limité, stop & simple~; le courtier achemine les ordres (livre~1, chapitre~10) \\
algorithme d'exécution & \omterm{def:mx:the-limit-order-book:orders}{ordres à cours limité} au meilleur cours, \omterm{def:mx:order-types-and-their-uses:peg}{ordres indexés}, cachés, iceberg, IOC vers des dark pools & travailler un ordre de grande taille sans l'afficher~; ne prendre de la liquidité que lorsque c'est nécessaire \\
teneur de marché & \omterm{def:mx:the-limit-order-book:orders}{ordres à cours limité} post-only, cotations en masse, IOC pour se couvrir & gagner le spread et le rabais~; ne jamais prendre par accident \\
trader sensible à la latence & IOC et ordres de balayage intermarchés & prendre un cours périmé avant qu'il ne bouge, ne toucher à rien d'autre \\
trader protégeant sa position & \omterm{def:mx:order-types-and-their-uses:stop}{ordres stop} et \omterm{def:mx:order-types-and-their-uses:stop}{ordres stop à cours limité} & sortir à un cours fixé sans surveiller l'écran \\
\bottomrule
\end{tabular}
\end{center}

Les types d'ordres se multiplient parce que chacun résout le problème d'un participant dans le cadre des règles de priorité, et chacun peut changer qui négocie en premier. Le
cas du hook de 2015 est l'extrême~: des types d'ordres qui réévaluaient le cours d'un ordre pour éviter de bloquer celui d'une autre plateforme, avec une priorité en coulisse que
seuls certains membres connaissaient pleinement. La leçon pour quiconque négocie est de lire ensemble le règlement de marché et la spécification technique, et de tester le
comportement du moteur à l'aune des deux~; pour quiconque construit un moteur (chapitre~26), il faut écrire la règle en un seul endroit et en dériver la documentation.

\section{Un ordre à cours limité est une option gratuite}

\begin{definition}[Option gratuite de l'ordre à cours limité]\label{def:mx:order-types-and-their-uses:option}
L'\emph{option gratuite de l'ordre à cours limité}\index{option gratuite de l'ordre à cours limité} est le droit qu'un \omterm{def:mx:the-limit-order-book:orders}{ordre à cours limité} au repos accorde au reste du marché~: acheter à son
cours (une vente au repos est un call émis à ce prix d'exercice) ou vendre à son cours (un achat au repos est un put), exerçable par n'importe qui, à tout moment, jusqu'à ce
que l'ordre soit annulé.
\end{definition}

Copeland et Galai (1983) ont décrit de cette façon le coût de la fourniture de cotations~: les cours acheteur et vendeur d'un dealer sont un put et un call émis au profit de traders qui en savent davantage.
L'option est exercée lorsque la valeur franchit le cours avant que le détenteur puisse annuler. Une version européenne donne une borne inférieure à sa valeur.

\begin{proposition}[La valeur cédée par un ordre au repos]\label{prop:mx:order-types-and-their-uses:option}
Un ordre de vente repose à une distance $d\ge0$ au-dessus du prix efficient $P^\ast_0$, qui suit un mouvement brownien de volatilité $\sigma$ par racine carrée
de seconde. Si un trader qui connaît $P^\ast_T$ achète à cet ordre à l'instant $T$ chaque fois que $P^\ast_T$ dépasse son cours, l'ordre perd en moyenne
\[
C(d)=\sigma\sqrt T\,\varphi\!\left(\frac{d}{\sigma\sqrt T}\right)-d\left(1-\Phi\!\left(\frac{d}{\sigma\sqrt T}\right)\right).
\]
\end{proposition}

\begin{proof}
La perte vaut $(P^\ast_T-P^\ast_0-d)^+$ avec $P^\ast_T-P^\ast_0\sim\mathcal N(0,\sigma^2T)$~: il s'agit de la valeur du call de Bachelier (livre~2, chapitre~13).
\end{proof}

\omcode{code/microstructure/02-order-types-and-their-uses/python/mx_ordertypes.py}{202}{209}{L'option gratuite d'un ordre au repos.}\label{lst:mx:order-types-and-their-uses:option}

Pendant l'heure simulée de \texttt{firm.tape}, le prix efficient se déplace de 0.27 tick par racine carrée de seconde. Une vente située à un demi-tick au-dessus de la valeur qui reste pendant
10 secondes cède en moyenne 0.15 tick par action, et 0.60 tick si elle reste une minute~: la valeur d'un ordre se dégrade avec le temps pendant lequel il reste
sans surveillance, ce qui explique que les fournisseurs de liquidité annulent autant (chapitre~1). Le coût mesuré est plus élevé, car les informés peuvent exercer au meilleur
moment, et non à un moment fixe~: évalué au prix efficient 10 secondes après chaque transaction, le côté au repos des transactions contre des agresseurs informés a perdu 1.22
tick par action~; contre des agresseurs non informés, il a gagné 0.58~; le flux informé représentait 9.2~\% du volume, et l'ensemble des ordres au repos a gagné 0.41. Les chapitres~4 et~5
transforment cette arithmétique en théorie de la fourchette.

\section{Tutoriel~: un seul livre, tous les types d'ordres}\label{tut:mx:order-types-and-their-uses:tut}

\begin{tutorial}
\textbf{Objectif.} Observer les types d'ordres agir sur un seul moteur, provoquer une cascade de stops et chercher un iceberg dans le flux public. \textbf{État final~:}
le tableau de la section~2, \cref{fig:mx:order-types-and-their-uses:cascade,fig:mx:order-types-and-their-uses:iceberg}.
\begin{tutsteps}
\item \textbf{Validation.} Lire les règles de type d'ordre qui s'appliquent avant l'appariement~:
\omcode{code/firm/ordertypes/firm_ordertypes.py}{71}{92}{Validation d'un ordre saisi dans \texttt{firm.ordertypes}.}\label{lst:mx:order-types-and-their-uses:validate}
\item \textbf{Démonstration.} \texttt{mx\_ordertypes.showcase()} envoie les huit ordres à un moteur nu et renvoie les rapports de chacun.
\item \textbf{Cascade.} \texttt{stop\_cascade(n)} pour 0 à 20 stops, avec des \omterm{def:mx:order-types-and-their-uses:stop}{ordres stop} et avec des \omterm{def:mx:order-types-and-their-uses:stop}{ordres stop à cours limité}.
\item \textbf{Icebergs.} \texttt{iceberg\_study(mode, delay\_ms, jitter)} pour les quatre cas, avec le détecteur strict et le détecteur souple~; tracer avec
\texttt{fig\_ordertypes.py}.
\item \textbf{L'option.} \texttt{passive\_markout(session())} et \texttt{free\_option(0.5, sigma\_per\_sqrt\_second(session()), 10)}.
\end{tutsteps}
\textbf{Pour aller plus loin.} Espacer les stops de deux centimes et trouver l'espacement auquel la cascade s'arrête d'elle-même~; donner au détecteur le
cours de l'iceberg comme a priori (il connaît le niveau) et voir combien de fausses alertes cela supprime.
\end{tutorial}

\section{Construction~: la couche des types d'ordres}\label{bld:mx:order-types-and-their-uses:ordertypes}

\begin{build}
\bfield{Rôle} Un seul lieu pour les règles que chaque type d'ordre ajoute à l'appariement, utilisées par le moteur du simulateur de bourse dans trois langages.
\bfield{Interface} \texttt{EOrder} (un \texttt{firm\_lob.Order} avec propriétaire, firme, \omterm{def:mx:order-types-and-their-uses:tif}{durée de validité}, affichage, post-only, quantité affichée, réserve, quantité
minimale, groupe et mode de prévention de l'auto-appariement, prix de stop, limite, conteneur)~; \texttt{validate(msg, inst, phase, frozen, band)} renvoyant un motif
de rejet~; \texttt{slice\_for\_display}~; \texttt{peg\_target(o, book, pegged)}~; \texttt{stp\_conflict}, \texttt{stp\_actions(mode)}~;
\texttt{stop\_triggered(o, last)}~; \texttt{marketable(side, limit, price)}.
\bfield{Règles} Les motifs de validation sont les codes de rejet du protocole de saisie des ordres~; un \omterm{def:mx:order-types-and-their-uses:peg}{ordre indexé au mid} est fixé à partir des cotations affichées et plafonné par
sa limite~; un \omterm{def:mx:order-types-and-their-uses:peg}{ordre indexé} primaire suit le meilleur prix non indexé affiché de son côté~; la quantité minimale n'est admise que pour les ordres IOC et les \omterm{def:mx:order-types-and-their-uses:peg}{ordres indexés au mid}~; une quantité
affichée n'est admise que pour les ordres visibles~; les stops se déclenchent sur les transactions, non sur les cotations.
\bfield{Tests d'acceptation} \texttt{code/firm/exchsim/tests/test\_engine\_rules.py}~: post-only, IOC et FOK, \omterm{def:mx:the-limit-order-book:orders}{ordres au marché} et bandes, icebergs, prévention de
l'auto-appariement dans chaque mode, priorité au remplacement, \omterm{def:mx:order-types-and-their-uses:peg}{ordres indexés au mid} avec un minimum, stops, la phase d'ouverture, les limitations de débit et les doublons~; les moteurs C++20 et Rust
reproduisent le moteur Python octet pour octet sur le journal partagé (chapitre~26).
\bfield{Pour aller plus loin} Ordres discrétionnaires (affichés à un prix, prêts à être exécutés à un autre)~; un ordre à prix glissant qui se reprice au lieu de bloquer un
marché éloigné, avec sa priorité énoncée.
\end{build}

\begin{omsources}
\item SEC, \emph{In the Matter of EDGA Exchange, Inc. and EDGX Exchange, Inc.}, Release No.~74032, 12~janvier 2015.
\item Nasdaq Rule 4703 (Order Attributes), tel que reproduit dans SR-NASDAQ-2026-055, Release No.~34-105718, juin 2026.
\item T.~E.~Copeland et D.~Galai, «~Information effects on the bid-ask spread~», \emph{Journal of Finance} 38(5), 1983.
\end{omsources}

\section{Exercices}

\begin{exercise}[$\star$]\label{exo:mx:order-types-and-their-uses:1}
Le meilleur cours vendeur est 100.01 pour 300 actions et 100.02 pour 500. Que fait un ordre d'achat à cours limité à 100.02 pour 600 actions, et que laisse-t-il dans le carnet d'ordres~?
\end{exercise}

\begin{exercise}[$\star$]\label{exo:mx:order-types-and-their-uses:2}
Avec 300 actions offertes à 100.01 et rien derrière elles dans la limite, que font un \omterm{def:mx:order-types-and-their-uses:tif}{ordre tout-ou-rien} et un \omterm{def:mx:order-types-and-their-uses:tif}{ordre exécuter-ou-annuler} d'achat de 400 à 100.01 chacun~?
\end{exercise}

\begin{exercise}[$\star$]\label{exo:mx:order-types-and-their-uses:3}
Pourquoi un \omterm{def:mx:order-types-and-their-uses:hidden}{ordre caché} se classe-t-il derrière un ordre affiché au même prix, même s'il est arrivé en premier~?
\end{exercise}

\begin{exercise}[$\star\star$]\label{exo:mx:order-types-and-their-uses:4}
Un iceberg de 10\,000 actions en affiche 500. S'il est exécuté en totalité, combien de tranches apparaissent sur le flux, et combien de fois il repasse en
queue de la file~?
\end{exercise}

\begin{exercise}[$\star\star$]\label{exo:mx:order-types-and-their-uses:5}
Une cotation repose à un tick de la valeur~; le prix efficient se déplace de 0.5 tick par racine carrée de seconde. Que donne la proposition pour une cotation laissée
pendant 30 secondes~?
\end{exercise}

\begin{exercise}[$\star\star$]\label{exo:mx:order-types-and-their-uses:6}
Expliquer pourquoi les \omterm{def:mx:order-types-and-their-uses:stop}{ordres stop à cours limité} raccourcissent la cascade de la \cref{fig:mx:order-types-and-their-uses:cascade}, et ce que leurs détenteurs abandonnent.
\end{exercise}

\begin{exercise}[$\star\star\star$]\label{exo:mx:order-types-and-their-uses:7}
\emph{Programmation.} Faire choisir à l'agent iceberg, comme prix de recharge, le meilleur cours vendeur courant au lieu de son premier prix, relancer le détecteur strict et expliquer
le changement.
\end{exercise}

\begin{exercise}[$\star\star\star$]\label{exo:mx:order-types-and-their-uses:8}
\emph{Trouver la faille.} «~Les \omterm{def:mx:order-types-and-their-uses:peg}{ordres indexés au mid} économisent la moitié de la fourchette sur chaque transaction, donc un algorithme d'exécution devrait envoyer tous ses ordres sous forme d'\omterm{def:mx:order-types-and-their-uses:peg}{ordres indexés au mid}.~»
\end{exercise}

\section{Problème~: l'iceberg dans le flux}

\begin{problem}[{Problème du week-end --- l'iceberg dans le flux}]\label{pb:mx:order-types-and-their-uses:1}
Un fonds vend 30\,000 actions au cours vendeur, 300 affichées à la fois. Un concurrent surveille le flux public. On dispose de \texttt{firm.exchsim}, du flux de fond du livre~7
et du détecteur du chapitre.

\medskip\noindent\textbf{Partie I --- Les types d'ordres.}
\begin{enumerate}
\item Quelle est la différence entre un iceberg et un \omterm{def:mx:order-types-and-their-uses:hidden}{ordre caché} en matière de priorité et de ce que le flux montre~?
\item Que devient la priorité d'un iceberg lorsque sa tranche est rafraîchie, et pourquoi~?
\item Contre quoi un \omterm{def:mx:order-types-and-their-uses:tif}{ordre tout-ou-rien} protège-t-il son émetteur, et que coûte-t-il~?
\item Pourquoi les \omterm{def:mx:order-types-and-their-uses:stop}{ordres stop} font-ils chuter plus vite un carnet qui baisse~?
\end{enumerate}

\medskip\noindent\textbf{Partie II --- La cascade.}
\begin{enumerate}[resume]
\item De combien une vente au marché de 3\,000 actions fait-elle bouger le prix sans stops, et avec vingt~?
\item Combien d'actions s'échangent dans le cas des vingt stops~?
\item De combien le prix baisse-t-il avec des \omterm{def:mx:order-types-and-their-uses:stop}{ordres stop à cours limité} placés un centime sous leur déclenchement~?
\item Que modifieriez-vous dans l'échelle des cours acheteurs pour arrêter la cascade~?
\end{enumerate}

\medskip\noindent\textbf{Partie III --- Le détecteur.}
\begin{enumerate}[resume]
\item Décrire le détecteur strict.
\item Combien de recharges l'iceberg de la plateforme effectue-t-il, et combien le détecteur strict en repère-t-il~? Avec combien de fausses alertes~?
\item Pourquoi un algorithme qui recharge aussitôt est-il repéré moins souvent que l'iceberg de la plateforme~?
\item Qu'est-ce qu'un retard aléatoire de 200 millisecondes et une taille aléatoire de 30~\% changent à la détection~?
\end{enumerate}

\medskip\noindent\textbf{Partie IV --- La course.}
\begin{enumerate}[resume]
\item Que détecte le détecteur souple avec les recharges retardées et randomisées, et à quel coût en fausses alertes~?
\item Énoncer le \emph{résultat nommé}~: les taux de détection et de fausses alertes du détecteur de recharges en fonction du retard et de la randomisation de la taille des recharges.
\item Quelle information supplémentaire le surveillant pourrait-il utiliser pour séparer les vraies recharges des coïncidences~?
\item Que perd le vendeur en retardant ses recharges~?
\item Est-il légal de détecter des icebergs à partir de données publiques~? Et de négocier sur la détection~?
\item Quel type d'ordre utiliseriez-vous pour vendre 30\,000 actions sans laisser apparaître de schéma de recharge~?
\item Comment une plateforme concevrait-elle un iceberg plus difficile à détecter~?
\item En une phrase~: que révèle un type d'ordre sur son propriétaire~?
\end{enumerate}
\end{problem}

\section{Questions d'entretien}

\begin{interviewq}[$\star$ \iqroles{trader}]\label{iq:mx:order-types-and-their-uses:1}
Quand utiliseriez-vous un \omterm{def:mx:order-types-and-their-uses:tif}{ordre exécuter-ou-annuler} plutôt qu'un \omterm{def:mx:the-limit-order-book:orders}{ordre à cours limité} qui reste au carnet~?
\end{interviewq}

\begin{interviewq}[$\star$ \iqroles{trader, chercheur}]\label{iq:mx:order-types-and-their-uses:2}
Qu'est-ce qu'un \omterm{def:mx:order-types-and-their-uses:stop}{ordre stop} de limitation de perte, et pourquoi peut-il s'exécuter loin de son prix de stop~?
\end{interviewq}

\begin{interviewq}[$\star\star$ \iqroles{chercheur}]\label{iq:mx:order-types-and-their-uses:3}
Pourquoi un \omterm{def:mx:the-limit-order-book:orders}{ordre à cours limité} au repos ressemble-t-il à une option vendue~? Qui détient l'option, et quel est son prix d'exercice~?
\end{interviewq}

\begin{interviewq}[$\star\star$ \iqroles{développeur}]\label{iq:mx:order-types-and-their-uses:4}
Comment implémenteriez-vous un \omterm{def:mx:order-types-and-their-uses:hidden}{ordre iceberg} dans un \omterm{def:mx:the-limit-order-book:engine}{moteur d'appariement}~? Qu'est-ce qui apparaît sur le flux lorsque la tranche affichée est exécutée~?
\end{interviewq}

\begin{interviewq}[$\star\star$ \iqroles{trader, développeur}]\label{iq:mx:order-types-and-their-uses:5}
Un \omterm{def:mx:order-types-and-their-uses:peg}{ordre indexé au mid} est au repos lorsque le marché devient bloqué (cours acheteur égal au cours vendeur). Que doit faire le moteur, et pourquoi~?
\end{interviewq}

\begin{interviewq}[$\star\star\star$ \iqroles{chercheur, trader}]\label{iq:mx:order-types-and-their-uses:6}
Comment détecteriez-vous un grand vendeur caché à partir des seules données de marché publiques~?
\end{interviewq}
